\subsection{Two further vanishing degrees and their character consequences}
\label{sec:k28-k36-vanishing}

\begin{theorem}[K17-K28-K36]\label{thm:k28-k36-vanish}
The stipulated arithmetic theorem of Alp\"oge implies
\[
                         K_{28}(\mathbf Z)=K_{36}(\mathbf Z)=0.
 \tag{HK2.1}
\]
Together with the preceding results, the established positive degrees
of the form $4k$ now include
\[
             K_{4k}(\mathbf Z)=0
             \qquad\bigl(k\in\{1,2,3,4,5,6,7,9\}\bigr).
 \tag{HK2.2}
\]
For every odd prime $\ell$, the cyclotomic class-group components
$A_\ell^{[-14]}$ and $A_\ell^{[-18]}$ vanish, and the strict groups
$\mathcal S_\ell(14)$ and $\mathcal S_\ell(18)$ vanish.
The standard full-domain Soul\'e characters
\[
 \kappa_{15}:N_\ell\longrightarrow\mathbf Z_\ell(15),
 \qquad
 \kappa_{19}:N_\ell\longrightarrow\mathbf Z_\ell(19),
 \quad
 N_\ell=\operatorname{Gal}
  (\overline{\mathbf Q}/\mathbf Q(\mu_{\ell^\infty})),
\]
are surjective. Consequently their reductions are surjective modulo
every positive power of $\ell$.
\end{theorem}

\begin{proof}
Set $e=14$ or $e=18$.
For every prime $\ell>2^{31}$, use the same closed source datum
$\mathrm{GD}(\Omega,E,\ell,e,1)$ with
$E=\mathbf Q(\sqrt{-3})$ and modulus two.
Its hypotheses hold: $e$ is positive and even, $\ell\ne2,3$,
and $\ell-1>2e$.
Corollary~\ref{cor:u3-general-coverage} places every actual admitted
rank-three system in the full compact coefficient space of highest
weight $(e-1,0,1-e)$.
Lemma~\ref{lem:gl2-conductor} places every admitted rank-two system
in the full classical space $S_{2e+1}(\Gamma_0(3),\chi_{-3})$.
These are exactly the spaces, scales, and trace normalizations used in
Proposition~\ref{prop:hecke-higher2}.

For $e=14$, the exact first row of (HH.2) has all prime factors
less than $2^{31}$.
Therefore, at the given $\ell$, at least one of $N_{14,7}$ and
$N_{14,13}$ is an $\ell$-adic unit.
Choose the corresponding auxiliary prime $p=7$ or $13$.
It is split in $E$ and automatically differs from $\ell$.
Corollary~\ref{cor:general-hecke-unit} makes the source strict
negative-character Selmer length zero.
The exact bridge (K5) gives
$K_{28}(\mathbf Z)\{\ell\}=0$.

For $e=18$, the same argument applies at every prime
$\ell>2^{31}$ except possibly
\[
                              q=8361790201,
\]
the sole large factor in the third row of (HH.2).
Proposition~\ref{prop:bernoulli-residue36}, including its two independent
complete integer calculations, gives
\[
                B_{q-19}\equiv1827838999\not\equiv0\pmod q.
\]
The index is even and lies between $2$ and $q-3$;
its Bernoulli denominator is prime to $q$.
Apply Lemma~\ref{lem:isolated-reflection} with $e=18$.
Herbrand's implication first gives $A_q^{[19]}=0$,
and the proved even-to-mirror reflection gives $A_q^{[-18]}=0$.
Then (K2) gives $K_{36}(\mathbf Z)\{q\}=0$.
This treats one specified component at $q$ and requires no
Vandiver assertion for its other characters.

For every odd prime $\ell<2^{31}$, the already cited published
Hart--Harvey--Ong Vandiver verification gives both primary vanishings
through (K2)--(K3). There is no prime equal to $2^{31}$.
The 2-primary parts vanish by the Rognes--Weibel theorem recorded
after (K3). Since these positive even $K$-groups are finite,
all primary parts vanish and (HK2.1) follows.
The preceding degree-twelve and higher-degree theorems, with the known
degrees four and eight, give (HK2.2).

The class-group assertions now follow from (K2), and the strict Selmer
assertions from (K5).
Finally apply the exact full-domain index formula (K18) at
$m=15$ and $m=19$:
\[
 [\mathbf Z_\ell(m):\kappa_m(N_\ell)]
       =|K_{2m-2}(\mathbf Z)\{\ell\}|=1.
\]
This proves the stated surjectivity and, by reduction, surjectivity
modulo every $\ell^a$.
The normalization is the cyclotomic product (K19); it does not insert
the universal-series factorial or restrict the domain to a deeper
Galois filtration subgroup.
\end{proof}

\begin{corollary}[K17-EIGHT-DEGREES-PERIOD]
\label{cor:eight-degrees-period}
For every odd prime $\ell$ and every positive integer $k$ whose residue
modulo $(\ell-1)/2$ is represented by one of
\[
                              0,1,2,3,4,5,6,7,9,
\]
one has $K_{4k}(\mathbf Z)\{\ell\}=0$.
\end{corollary}

\begin{proof}
For a nonzero listed representative use (HK2.2) and (K2):
the relevant exponent $-2k$ is unchanged modulo $\ell-1$ when $k$
is changed by a multiple of $(\ell-1)/2$.
The residue zero gives the trivial class-group component, which
vanishes by extension and norm to $\mathbf Q$, as proved after (K3).
The assertion concerns vanishing, not an isomorphism or periodicity
of the full finite groups.
\end{proof}

\paragraph{The middle row before changing the quadratic datum.}
The $e=16$ comparison over this quadratic field retains the isolated prime
$25078104767246023$.
This pair of comparisons alone does not determine that component.
An explicit sufficient Bernoulli condition would be
$B_{25078104767246006}\not\equiv0
 \pmod{25078104767246023}$, as specified by
Lemma~\ref{lem:isolated-reflection}. Theorem~\ref{thm:k32-vanish}
below instead closes the component by a different quadratic datum
and proves $K_{32}(\mathbf Z)=0$, without evaluating that Bernoulli
residue.
